@@M0@@; we have the commutative diagram:
@@M1@@
where the morphisms @@M2@@ and @@M3@@ are surjective. Let @@M4@@, for @@M5@@, be elements of @@M6@@ such that for every @@M7@@ we have
@@M8@@
One defines a homomorphism @@M9@@ by setting @@M10@@. This gives a closed immersion @@M11@@ which is a section of @@M12@@, and which satisfies @@M13@@. Moreover, being a section of a smooth morphism, it is a regular immersion (VII 1.10). Knowing that @@M14@@ and @@M15@@ are regular, we shall deduce that @@M16@@ is regular if we show that the sequence:
@@M17@@
is exact and locally split, where @@M18@@, @@M19@@ and @@M20@@ denote respectively the ideals defining @@M21@@, @@M22@@, and @@M23@@ (VII 1.7). For this it suffices to show that @@M24@@ is a direct summand of @@M25@@. Now one sees at once that @@M26@@ is obtained from @@M27@@ by the base change @@M28@@. Furthermore, the immersions @@M29@@ and @@M30@@ are regular (VII 1.10). Consequently (VII 1.2), the ideal defining @@M31@@ is @@M32@@, and the conormal sheaf of @@M33@@ is @@M34@@. The exact sequence of conormal sheaves for the decomposition @@M35@@ therefore gives, in particular, a homomorphism: @@M36@@. It is immediate to check that the composite: @@M37@@ is the identity (from the affine point of view, this corresponds to the decomposition of @@M38@@ as the direct sum of an ideal of @@M39@@, and the ideal generated by the @@M40@@). This indeed gives the required splitting.
